% ch10-ranking — ranking objectives: who gets funded next (target T1)

\section{The metric is the ranking}\label{ch10:sec:target}

\Cref{ch:question} argued that the first answer to ``what is valuable to
know'' is an identity, not a time: given finite diligence capacity, which
$K$ firms should be looked at now because they are most likely to be the
next to raise. The consuming action is a short list, refreshed on a
cadence; nothing in that action consumes an arrival time. The target is
therefore the conditional distribution of the next deal's identity,
already fixed as \cref{eq:ch01:t1},
\begin{equation}\label{eq:ch10:target}
p(i \st s, t)
 \;=\;
 \Prob\big(\text{next deal in sector } s \text{ goes to firm } i
 \;\big|\; \text{a deal occurs},\; \Ftm\big),
 \qquad i \in \riskset_s(t),
\end{equation}
where $\riskset_s(t) \subseteq \firms$ is the risk set of tracked firms at
risk in sector $s$ at $t^-$. This chapter's thesis is a sharpening of that
framing: when T1 is the valuable thing, \emph{the likelihood of the full
marked process is not the metric --- the ranking is}. The full likelihood
rewards timing accuracy that no action consumes and that the data defects
of \cref{ch:data} corrupt; the object to be scored is the ordered list,
and the object to be fitted is the conditional law
\cref{eq:ch10:target}, whose own likelihood --- the conditional NLL of the
mark factor in \cref{thm:ch03:groundmark} --- is the proper training
loss precisely because it is a likelihood \emph{of the ranking target and
of nothing else}.

The bridge between intensity models and \cref{eq:ch10:target} is the
competing-intensities identity. \Cref{ch:foundations} states it as part of
the ground--mark factorization; here we give the operational proof,
because every robustness claim in this chapter is a corollary of the
\emph{form} of the identity --- a ratio.

\begin{assumption}[Firm-level intensity system]\label{asm:ch10:intensities}
Fix a sector $s$ and work conditionally on $\Ftm$. Each firm $i \in
\riskset_s(t)$ has a counting process $N_i$ with $\F$-conditional
intensity $\lam_i(t) \ge 0$ (\cref{def:ch02:intensity}), the intensities
are locally integrable, no two firms jump at the same instant almost
surely, and the sector ground process $N_s = \sum_{i} N_i$ has intensity
$\lam_{\mathrm g}(t) = \sum_{j \in \riskset_s(t)} \lam_j(t) \in
(0,\infty)$.
\end{assumption}

\begin{proposition}[Competing intensities]\label{prop:ch10:competing}
Under \cref{asm:ch10:intensities}:
\begin{enumerate}[label=(\roman*)]
\item \emph{(Exact form, at the event.)} For almost every $t$, the
conditional distribution of the identity of an event occurring at $t$ is
\begin{equation}\label{eq:ch10:competing}
\Prob\big(\text{identity} = i \st \text{event at } t,\; \Ftm\big)
 \;=\;
 \frac{\lam_i(t)}{\sum_{j \in \riskset_s(t)} \lam_j(t)} .
\end{equation}
\item \emph{(Next-event form.)} Fix a decision time $\tau$ and suppose in
addition that on the no-event continuation after $\tau$ the intensities
are $\F_\tau$-measurable functions $u \mapsto \lam_i^{\tau}(u)$ (true for
internal-history models with exogenous covariate paths frozen at $\tau$).
Then the next event time $T$ has conditional density
$\lam^{\tau}_{\mathrm g}(u)\exp(-\int_\tau^u \lam^{\tau}_{\mathrm g})$
and, jointly with the identity $J$,
\begin{equation}\label{eq:ch10:nextevent}
\Prob\big(J = i,\; T \in \dd u \st \F_\tau\big)
 \;=\;
 \lam_i^{\tau}(u)\,
 \exp\!\Big(-\!\int_\tau^{u} \lam^{\tau}_{\mathrm g}(v)\dd v\Big) \dd u,
\qquad
\Prob\big(J = i \st T = u, \F_\tau\big)
 \;=\;
 \frac{\lam_i^{\tau}(u)}{\lam^{\tau}_{\mathrm g}(u)} .
\end{equation}
\end{enumerate}
\end{proposition}

\begin{proof}
(i) View the system as a single marked point process on $[0,\Tend] \times
\riskset_s(t)$ with intensity kernel $\lam(t, \{i\}) = \lam_i(t)$; the
absence of common jumps makes this kernel well defined, and superposition
gives the ground intensity $\lam_{\mathrm g} = \sum_j \lam_j$. The
ground--mark factorization of \cref{thm:ch03:groundmark} writes the
kernel as $\lam(t,\{i\}) = \lam_{\mathrm g}(t)\, p(i \st t, \Ftm)$, where
$p(\cdot \st t, \Ftm)$ \emph{is}, by construction, the conditional mark
distribution given an event at $t$ \citep[Ch.~6]{daley2003}. Identifying
the two expressions for the kernel and dividing by $\lam_{\mathrm g}(t) >
0$ yields \cref{eq:ch10:competing}.

(ii) The operational derivation. Partition $(\tau, u]$ into $n$ intervals
of mesh $\delta = (u-\tau)/n$. Given $\F_\tau$ and the frozen intensity
functions, the probability that no firm jumps in $(v, v+\delta]$, given no
jump before $v$, is $1 - \lam^{\tau}_{\mathrm g}(v)\delta + o(\delta)$,
and these no-jump factors multiply across subintervals because, in the
absence of events, the conditional jump probabilities depend only on the
deterministic functions $\lam^{\tau}_j$ --- this is the
independent-increment approximation, exact in the limit $\delta \to 0$ by
the product-integral representation of the survivor function
\citep{abgk1993}. Hence $\Prob(T > u \st \F_\tau) = \exp(-\int_\tau^u
\lam^{\tau}_{\mathrm g})$. The event $\{T \in (u, u+\delta],\, J = i\}$
requires no jump on $(\tau, u]$ and a jump of firm $i$ in $(u,
u+\delta]$, of probability $\lam^{\tau}_i(u)\delta + o(\delta)$;
multiplying and letting $\delta \to 0$ gives the joint density in
\cref{eq:ch10:nextevent}. Dividing by the marginal density of $T$
--- the same expression with $\lam^{\tau}_i$ replaced by
$\lam^{\tau}_{\mathrm g}$ --- gives the ratio. When exogenous covariates
evolve stochastically after $\tau$, part (ii) holds with the intensities
replaced by their conditional expectations along the no-event
continuation; part (i) is unaffected, and as $u \downarrow \tau$ the two
parts agree, so ranking \emph{at the decision time} inherits the exact
form (i).
\end{proof}

\begin{corollary}[Every intensity model is a ranker]\label{cor:ch10:induced}
Any firm-level intensity model $\{\lam_i\}$ induces a ranker: order
$\riskset_s(t)$ by $\lam_i(t)$, equivalently by the conditional
probability \cref{eq:ch10:competing}, of which $\lam_i(t)$ is a monotone
transform. Under the model, this ordering is optimal for every list
metric considered in this chapter (\cref{lem:ch10:topk}). Conversely ---
and this is the direction the chapter exploits --- the ratio
\cref{eq:ch10:competing} can be modeled \emph{directly}, without ever
committing to a timing model for $\lam_{\mathrm g}$.
\end{corollary}

The ratio form is the load-bearing observation. Any factor common to the
whole risk set --- the sector's level of activity, the calendar week's
deal flow, a macro regime, an unmodeled burst --- appears in every
$\lam_j$ and cancels. The next section turns this cancellation into a
likelihood with a proof of exactly what it buys.

\section{The conditional-choice mark factor}\label{ch10:sec:choice}

The repository's mark factor models the ratio directly. Each candidate
$i \in \riskset_s(t)$ receives a score built from point-in-time features
$\zfeat_i(t)$ (firm covariates with staleness ages, \cref{ch:stale};
engineered history summaries, \cref{def:ch02:covariates}) and a cooldown
indicator $C_i(t) = \ind{\text{firm } i \text{ funded in the last } K_c
\text{ weeks}}$:
\begin{equation}\label{eq:ch10:score}
\score_i(t)
 \;=\;
 (w_0 + u_s)\T \zfeat_i(t) \;+\; \eta_s\, C_i(t),
 \qquad \eta_s \le 0,
\end{equation}
with a global weight vector $w_0$, sector deviations $u_s$, and
sector-specific cooldown coefficients $\eta_s$. The conditional choice
probability is the softmax over the \emph{current} risk set,
\begin{equation}\label{eq:ch10:softmax}
p(i \st s, t)
 \;=\;
 \frac{\exp \score_i(t)}{\sum_{j \in \riskset_s(t)} \exp \score_j(t)},
\end{equation}
and the training loss is the conditional negative log-likelihood over
observed events $(t_m, s_m, i_m)$,
\begin{equation}\label{eq:ch10:nll}
-\loglik_{\mathrm c}(\thetamodel)
 \;=\;
 \sum_{m}
 \Big[
 \logsumexp_{j \in \riskset_{s_m}(t_m)} \score_j(t_m)
 \;-\;
 \score_{i_m}(t_m)
 \Big].
\end{equation}
The normalization over $\riskset_s(t)$ rather than over a fixed candidate
list is the crucial design choice: firms enter the risk set by becoming
active and leave it by the exposure rules of \cref{ch:buckets}, and the
model is well defined for every risk-set size. Algebraically,
\cref{eq:ch10:softmax,eq:ch10:nll} are McFadden's conditional logit
\citep{mcfadden1974,train2009}; they are simultaneously, term for term,
the factors of Cox's partial likelihood for a proportional-intensity
model \citep{cox1972,cox1975}. The two literatures prove the same theorem
twice, and we will lean on both.

\begin{proposition}[Baseline cancellation]\label{prop:ch10:cancel}
Suppose the firm intensities have the multiplicative form
\begin{equation}\label{eq:ch10:propint}
\lam_i(t) \;=\; \lam_0(s, t)\, \exp \score_i(t),
\qquad i \in \riskset_s(t),
\end{equation}
where $\lam_0(s,t) \ge 0$ is \emph{any} sector--time baseline common to
the risk set: unknown, unmodeled, nonstationary, dependent on the whole
history, or misspecified in any way whatsoever. Then the conditional
identity distribution given an event is \cref{eq:ch10:softmax}, free of
$\lam_0$. Consequently the conditional likelihood \cref{eq:ch10:nll}
involves no timing model, and any deformation of the time axis or of the
baseline that acts on the risk set as a common factor --- weekly
aggregation, jittered or back-filled dates that preserve the assignment
of events to risk-set snapshots, bursts, regime shifts, unmodeled
excitation in the ground process --- leaves \cref{eq:ch10:nll} unchanged.
\end{proposition}

\begin{proof}
Substitute \cref{eq:ch10:propint} into \cref{eq:ch10:competing}:
\[
\frac{\lam_i(t)}{\sum_{j} \lam_j(t)}
 =
\frac{\lam_0(s,t) e^{\score_i(t)}}{\lam_0(s,t) \sum_{j} e^{\score_j(t)}}
 =
\frac{e^{\score_i(t)}}{\sum_{j \in \riskset_s(t)} e^{\score_j(t)}} .
\]
The baseline cancels because it is common; nothing about its form was
used. For the aggregation claim: replacing $t_m$ by its containing bin
changes which snapshot $(\riskset_s, \zfeat, C)$ the event is matched to,
but if features and risk sets are constructed at bin resolution from
pre-bin information (\cref{prot:ch02:pit}), the conditional law of the
winner given the bin's snapshot is again \cref{eq:ch10:softmax} with the
bin-level baseline cancelled --- interval censoring corrupts $\lam_0$,
which does not appear. What does \emph{not} cancel is instructive: a
firm-specific multiplicative frailty $\lam_i = \lam_0 \xi_i e^{\score_i}$
does not cancel (it is not common) and is absorbed into the score as an
omitted variable; and a baseline that differs \emph{within} the risk set
--- e.g.\ stage-specific deal clocks inside one sector --- breaks the
common-factor premise and must be moved into $\score_i$ as a covariate or
handled by stratifying the risk set.
\end{proof}

\begin{proposition}[Robustness of the conditional MLE]\label{prop:ch10:consistency}
Let events $m = 1, \dots, M$ carry risk sets and features constructed
point-in-time, with $i_m \in \riskset_{s_m}(t_m)$. Suppose (a) the
conditional model is correctly specified: there exists $\thetamodel^*$
such that $\Prob(i_m = i \st \text{event } m, \F_{t_m^-}) = p(i \st s_m,
t_m; \thetamodel^*)$ of \cref{eq:ch10:softmax}; (b) features are
uniformly bounded, $\thetamodel^*$ is identified and interior to a
compact parameter set; (c) the averaged conditional information
$\bar\Fisher_M = M^{-1}\sum_m \E[V_m]$, with $V_m$ the conditional
covariance of the score gradient under \cref{eq:ch10:softmax}, converges
to a positive-definite limit. Then the maximizer $\thetahat$ of
$\loglik_{\mathrm c}$ is consistent and asymptotically normal,
\emph{with no assumption on the ground process}: the law of $(t_m, s_m)$
may be arbitrarily misspecified, dependent, aggregated, or simply never
modeled.
\end{proposition}

\begin{proof}
Let $\mathcal{G}_m = \F_{t_m}$ be the event-indexed filtration. The
$m$-th score contribution at $\thetamodel^*$ is
\[
U_m \;=\; \nabla \score_{i_m}(t_m)
 \;-\; \sum_{j \in \riskset_{s_m}(t_m)} p(j \st s_m, t_m;\thetamodel^*)
 \, \nabla \score_j(t_m),
\]
and by (a), $\E[U_m \st \mathcal{G}_{m-1}, \text{event } m \text{ at }
(t_m, s_m)] = \sum_i p^*_i \nabla \score_i - \sum_j p^*_j \nabla \score_j
= 0$: the score is a martingale difference array with respect to
$\{\mathcal{G}_m\}$ \emph{conditionally on the event stream, whatever its
law}. Boundedness (b) and nondegeneracy (c) then give a law of large
numbers and a martingale central limit theorem for the score, and
consistency plus asymptotic normality of $\thetahat$ follow by the
standard argument --- this is precisely Cox's partial-likelihood
heuristic \citep{cox1975} made rigorous by counting-process martingale
theory in \citet{andersengill1982}; the conditional-logit sampling theory
of \citet{mcfadden1974} gives the same conclusion when events are treated
as independent choice occasions, and the martingale version shows
independence is not needed. The proof fails exactly where it should: if
the risk set is wrong (defect (C3): members retained who cannot win, or
competitors omitted), the softmax in $U_m$ is normalized over the wrong
set and the conditional mean of the score is no longer zero at
$\thetamodel^*$ --- the estimand silently changes to the conditional law
\emph{given the contaminated risk set}. Robustness to timing is not
robustness to membership.
\end{proof}

\Cref{prop:ch10:cancel,prop:ch10:consistency} together are the precise
version of the claim previewed in \cref{ch01:sec:factorization}: the mark
factor needs no timing model, survives the (C4) aggregation and
untrusted-date defect wholesale, and is insulated from ground-process
misspecification --- the three failure modes that \cref{ch:censoring}
shows corrupt every timing estimate on current data. That is why T1 is
the most defensible deliverable on present data quality, and why the
ranking layer, not the Hawkes layer, is the repository's headline
deliverable while the timing stack matures.

\section{Positive-only labels: manufacturing the comparison}\label{ch10:sec:positive}

A supervised-classification framing of T1 needs a negative class, and the
data has none (\cref{ch:data}): we observe fundings, never rejections. No
row anywhere says ``firm $i$ sought capital at $t$ and was declined,'' so
per-firm labels are positive-only, and any attempt to manufacture
firm-level negatives by declaring non-events (``firm $i$ did not raise in
week $t$'') imports the timing model through the back door --- the
probability of a non-event depends on $\lam_0$, which we just proved we
do not need and cannot trust. The conditional-choice construction solves
the label problem structurally: each funding event manufactures a valid
comparison --- the funded firm against its risk-set peers at the same
sector-time --- and \cref{prop:ch10:consistency} says these comparisons
identify the score parameters without a single explicit negative. This is
the same device by which Cox regression learns hazard ratios from cases
and risk sets alone, and the connection to case-control methodology is
exact: the risk set is the control pool, sampled at the case's own event
time.

The practical obstacle is scale: real sector risk sets contain hundreds
to thousands of firms, and \cref{eq:ch10:nll} sums over all of them per
event. Sampling the risk set fixes this without changing the estimand.

\begin{theorem}[Consistency of the sampled softmax]\label{thm:ch10:sampled}
For event $m$, let a subset $D_m \subseteq \riskset_{s_m}(t_m)$ with
$i_m \in D_m$ be drawn by a sampling protocol $\pi(D \st i, \riskset)$
--- the probability of drawing $D$ given that the winner is $i$ ---
satisfying the \emph{uniform conditioning property}: $\pi(D \st i) =
\pi(D \st j) > 0$ for all $i, j \in D$. Then, under the model
\cref{eq:ch10:softmax}, the conditional law of the winner given $D_m$ is
the softmax restricted to $D_m$,
\begin{equation}\label{eq:ch10:sampled}
\Prob\big(i_m = i \st D_m, \F_{t_m^-}\big)
 \;=\;
 \frac{\exp \score_i(t_m)}{\sum_{j \in D_m} \exp \score_j(t_m)},
\end{equation}
so maximizing the subsampled conditional likelihood built from
\cref{eq:ch10:sampled} is maximum likelihood in a correctly specified
model and inherits the consistency of \cref{prop:ch10:consistency}. In
particular, uniform sampling of $J$ non-winners from
$\riskset \setminus \{i_m\}$, with the winner always retained, satisfies
the property, and no correction term is needed. For a general protocol,
replacing $\score_j$ by $\score_j + \log \pi(D \st j)$ restores
\cref{eq:ch10:sampled} \citep{mcfadden1978,train2009}.
\end{theorem}

\begin{proof}
By Bayes' rule on the joint draw (winner first, then the sampled set
given the winner),
\[
\Prob(i_m = i \st D_m = D)
 \;=\;
 \frac{p(i \st s_m, t_m)\, \pi(D \st i)}
      {\sum_{j \in D} p(j \st s_m, t_m)\, \pi(D \st j)}
 \;=\;
 \frac{e^{\score_i + \log \pi(D \mid i)}}
      {\sum_{j \in D} e^{\score_j + \log \pi(D \mid j)}},
\]
where the full-risk-set normalizer of $p(\cdot \st s_m, t_m)$ cancels
between numerator and denominator --- the same ratio trick as
\cref{prop:ch10:cancel}, now applied to the sampling design. Under
uniform conditioning the $\log \pi$ terms are equal across $j \in D$ and
cancel as a common shift of the scores, giving \cref{eq:ch10:sampled}.
Uniform sampling of $J$ alternatives yields $\pi(D \st i) =
\binom{|\riskset| - 1}{J}^{-1}$ for every $i \in D$, which is constant.
The subsampled events $(i_m, D_m)$ then follow a correctly specified
conditional logit over $D_m$, and \cref{prop:ch10:consistency} applies
verbatim with $\riskset$ replaced by $D$. What is lost is efficiency, not
consistency: the conditional information per event shrinks with $J$, so
$J$ is a compute--variance dial, not a bias dial.
\end{proof}

\section{Cooldown: inhibition carried by the mark factor}\label{ch10:sec:cooldown}

\Cref{ch:hawkes} establishes that a nonnegative Hawkes kernel cannot
express inhibition and enumerates the three legal routes: (a) rectified
additive intensities, (b) exp-link log-intensity models, and (c)
inhibition carried by the mark factor. The cooldown term of
\cref{eq:ch10:score} is route (c). A firm that has just raised is, for a
while, less likely to raise again --- not because the sector cooled, but
because that firm left the market for capital. Encoding this as a
negative-coefficient covariate in the score gives, by
\cref{prop:ch10:cancel}, a multiplicative suppression of the firm's
\emph{share}: while $C_i = 1$, firm $i$'s conditional probability is
scaled by $e^{\eta_s} < 1$ relative to its own baseline share,
renormalized within the risk set. The ground layer is untouched --- sector
excitation stays positive and the linear Hawkes/MBPP mean-field structure
of \cref{ch:hawkes,ch:censoring} survives, which is exactly why route (c)
was chosen as the default (the positivity link and the expectation
operator do not commute for routes (a) and (b), breaking the closed-form
machinery).

The economic interpretation is direct: $\eta_s$ is the log of the factor
by which a recent raise suppresses a firm's odds of being the next funded
firm in its sector, and the constraint $\eta_s \le 0$ (enforced as a
bound in \code{fit\_startup\_ranker}) encodes the sign as a modeling
assumption rather than a hope. The rehearsal evidence says the device
works even under misspecification: in regime B --- where the generator's
inhibition is a continuous 30-day-half-life kernel, not a binary window
--- the fitted mark-factor cooldown recovered the generator's
self-inhibition \emph{ordering} across firms with rank correlation
$0.937$ (\cref{app:results}). The binary window $C_i(t)$ with default
$K_c = 26$ weeks is, of course, a one-column basis for a decaying curve;
the natural refinements --- an age-since-last-deal feature with learned
shape, or a few staggered windows --- belong to the gradient-boosted
scorer of the next section, which can learn the shape without a
parametric commitment.

\section{Gradient-boosted scores in the same likelihood}\label{ch10:sec:boosted}

The linear score \cref{eq:ch10:score} is a restriction, and it has been
falsified: the B1 rehearsal test fired, with a feature-parity XGBoost
ranker beating the linear conditional logit by $+0.22$ nats of held-out
conditional NLL and $+9$ points of top-5 hit rate on regime A
(\cref{app:results}). The design response is not to abandon the
conditional-choice likelihood but to change the score family inside it:
replace $\score_i(t)$ by $f(\zfeat_i(t), C_i(t), \staleage_i(t), \dots)$
with $f$ an additive tree ensemble \citep{friedman2001,chenguestrin2016},
keeping the softmax over the current risk set and the loss
\cref{eq:ch10:nll} exactly as they are. Training requires only the
per-candidate gradient and Hessian of the grouped-softmax loss, which is
where the ``listwise'' machinery of the learning-to-rank literature
collapses into three lines of calculus.

\begin{lemma}[Grouped-softmax gradients]\label{lem:ch10:gradhess}
Fix one event group $m$ with candidate scores $\{\score_j\}_{j \in
\riskset_m}$, softmax probabilities $p_j$, and one-hot label $y_j =
\ind{j = i_m}$. The per-event loss $\ell_m = \logsumexp_j \score_j -
\score_{i_m}$ has
\begin{equation}\label{eq:ch10:gradhess}
\frac{\partial \ell_m}{\partial \score_j} \;=\; p_j - y_j,
\qquad
\frac{\partial^2 \ell_m}{\partial \score_j \partial \score_k}
 \;=\; p_j \ind{j = k} - p_j p_k,
\end{equation}
so the diagonal Hessian entries are $h_j = p_j(1 - p_j)$. The full
Hessian $\diag(p) - p p\T$ is positive semidefinite (it is the
covariance matrix of the candidate-indicator vector under $p$), so
$\ell_m$ is convex in the scores and the diagonal is a legitimate
curvature surrogate.
\end{lemma}

\begin{proof}
Differentiate $\logsumexp$: $\partial_{\score_j} \logsumexp = p_j$;
subtracting $\partial_{\score_j} \score_{i_m} = y_j$ gives the gradient.
Differentiating $p_j$ once more gives $\partial_{\score_k} p_j = p_j
\ind{j=k} - p_j p_k$. For a vector $v$, $v\T(\diag(p) - pp\T)v = \E_p[v_J^2]
- (\E_p[v_J])^2 = \Var_p(v_J) \ge 0$.
\end{proof}

Implementation is a custom objective for XGBoost: candidates of one event
form one group; per boosting round, compute the softmax within each group
from the current margin scores, emit $g_j = p_j - y_j$ and $h_j = p_j(1 -
p_j)$ per candidate row, and let the tree learner take the regularized
Newton step in function space \citep{chenguestrin2016}. Dropping the
off-diagonal Hessian terms is the standard diagonal approximation; by
convexity of $\ell_m$ it costs step quality, never correctness of the
stationary point. Note what this objective is and is not: it \emph{is}
the exact conditional NLL of \cref{eq:ch10:nll} --- the listwise
cross-entropy per event group --- so the fitted trees produce genuine
conditional probabilities through the softmax; it is \emph{not} a
pairwise surrogate in the RankNet family \citep{burges2005}, which we
treat separately in \cref{ch10:sec:direct}. Sampled groups per
\cref{thm:ch10:sampled} keep the objective correct at scale.

\begin{warning}[The robustness belongs to the family, not the score]\label{warn:ch10:family}
Every robustness statement of \cref{ch10:sec:choice} ---
baseline cancellation, timing-free consistency, survival of aggregation
--- is a property of the conditional-choice \emph{likelihood family}: any
scorer, linear or boosted, normalized over the current risk set inherits
them. None of them certify the \emph{linear} score, which the B1 verdict
falsified, and none of them certify a boosted scorer trained on a
different loss (a pairwise surrogate, a pointwise classifier): change the
likelihood family and the cancellation proofs no longer apply. When this
chapter says ``the ranker is robust to (C1) and (C4),'' it means the
grouped-softmax likelihood is; claims about any particular fitted score
function are empirical and live in \cref{app:results}.
\end{warning}

\section{Multiple deals per sector-week: Plackett--Luce}\label{ch10:sec:ties}

Weekly resolution creates ties: several deals land in the same
sector-week cell and share a risk-set snapshot. The current
implementation treats them as repeated independent softmax draws over the
same risk set. Survival analysts will recognize the situation instantly:
tied event times in a Cox partial likelihood, where the repeated-softmax
treatment is exactly Breslow's approximation \citep{breslow1974}, the
exact treatment is the sequential-conditioning likelihood, and Efron's
correction sits between them \citep{efron1977,therneau2000}. The
discrete-choice literature derives the exact object independently as the
Plackett--Luce model \citep{plackett1975,luce1959}.

The derivation is a corollary of \cref{prop:ch10:competing}. Let $d$
deals occur in cell $(s,t)$ with winners $i_1, \dots, i_d$ in sequence.
After the $k$-th deal, firm $i_k$ ceases to compete for the remaining
within-cell deals (it has just raised; within one week its continued
presence in the denominator is a modeling error, not a modeling choice).
Applying \cref{eq:ch10:competing} to each successive deal over the
shrinking risk set gives the sequential softmax without replacement:
\begin{equation}\label{eq:ch10:pl}
\Prob\big(i_1, \dots, i_d \st s, t, d \text{ deals}\big)
 \;=\;
 \prod_{k=1}^{d}
 \frac{\exp \score_{i_k}(t)}
      {\sum_{j \in \riskset_s(t) \setminus \{i_1, \dots, i_{k-1}\}}
       \exp \score_j(t)} ,
\end{equation}
which is the Plackett--Luce likelihood of the ordered top-$d$ under score
utilities. The repeated-softmax approximation replaces every denominator
by the full-risk-set sum, so its per-cell log-likelihood error is
\begin{equation}\label{eq:ch10:plbias}
\Delta \ell_{s,t}
 \;=\;
 \sum_{k=2}^{d} -\log\big(1 - P_{k-1}\big),
 \qquad
 P_{k-1} \;=\; \sum_{l < k} p(i_l \st s, t),
\end{equation}
the cumulative softmax mass of already-funded firms. Two consequences.
First, the sign is systematic: the approximation understates every
later-deal probability, so it penalizes exactly the concentrated,
confident score profiles the ranker is supposed to produce, biasing
fitted scores toward diffuseness. Second, the size is quantifiable: to
first order $\Delta \ell_{s,t} \approx \binom{d}{2} \bar p$, with $\bar
p$ the typical winner probability --- negligible when risk sets are large
and diffuse ($\bar p \sim 10^{-2}$, $d \le 2$), material when many deals
hit a small or concentrated risk set in one week ($d \ge 3$ with $\bar p
\sim 10^{-1}$ costs tenths of a nat per cell, the same order as the
entire B1 improvement). Since same-week multiplicity is common in hot
sectors --- precisely the cells the business cares about --- the exact
likelihood is worth its implementation cost. One honest complication:
\cref{eq:ch10:pl} conditions on the within-cell \emph{order}, and (C4)
says within-week order is exactly what the dates cannot support. The
clean treatment is the exchangeable-order PL: average the log-likelihood
over random permutations of the cell's winners (exact summation over
$d!$ orders for the typical $d \le 3$, sampled permutations beyond),
which removes the dependence on untrusted ordering at the price of a
small, controlled Monte Carlo cost.

\begin{decision}{D10.2 --- Plackett--Luce in the next ranker revision}
The next revision of the mark factor replaces repeated softmax with the
Plackett--Luce likelihood \cref{eq:ch10:pl}, with within-cell order
handled by exact averaging over permutations for $d \le 3$ and sampled
permutations beyond, and with cooldown state updated between within-cell
draws in simulation (as \code{simulate\_marked\_paths} already does).
Rationale: \cref{eq:ch10:plbias} gives a systematic, concentration-punishing
bias whose magnitude in hot cells reaches the scale of the B1 gain; the
Cox-ties literature says the exact method dominates Breslow whenever ties
are heavy \citep{therneau2000}. Reversal trigger: if on real data the
measured NLL and recall@$K$ differences between PL and repeated softmax
are inside the era-blocked bootstrap intervals of
\cref{prot:ch10:recall} for two consecutive evaluation campaigns, revert
to repeated softmax for simplicity and record the equivalence.
\end{decision}

\section{Optimizing the list directly}\label{ch10:sec:direct}

If the consumed object is only the top-$K$ list, one can ask whether the
likelihood should be bypassed entirely in favor of a loss that targets
the list metric --- the learning-to-rank program: pairwise RankNet
\citep{burges2005}, listwise LambdaMART with its $\lambda$-gradients
aimed at NDCG \citep{burges2010,jarvelin2002}. The temptation deserves a
precise answer, and the answer has two halves: a lemma that says
likelihood training is already list-optimal \emph{under the model}, and
an empirical clause that says surrogates can win only through
misspecification --- which B1 proved is real.

\begin{lemma}[Top-$K$ optimality of the conditional probabilities]\label{lem:ch10:topk}
Fix a sector-time $(s, \tau)$ and let $p_i$, $i \in \riskset_s(\tau)$, be
the conditional identity distribution of the next deal. For a list $L
\subseteq \riskset_s(\tau)$ with $|L| = K$, the expected recall@$K$
against the next deal is $\E[\ind{\text{winner} \in L}] = \sum_{i \in L}
p_i$, and it is maximized by any $L$ consisting of $K$ largest-$p_i$
candidates. For an ordered list $\sigma$, the expected reciprocal rank
$\sum_{r \ge 1} p_{\sigma(r)} / r$ is maximized by sorting $p$ in
descending order. Hence, under the model, producing the top-$K$ by
conditional probability is the optimal list policy for both headline
metrics, and NDCG with flat gains and any decreasing discount is
maximized by the same sort.
\end{lemma}

\begin{proof}
Recall: the objective $\sum_{i \in L} p_i$ is a set function maximized by
exchange: if $i \in L$, $i' \notin L$, and $p_{i'} > p_i$, swapping $i'$
for $i$ strictly increases the sum; iterating terminates at a set of $K$
largest probabilities. MRR: the discounts $1/r$ are strictly decreasing;
if an ordering places $p_{\sigma(r)} < p_{\sigma(r')}$ with $r < r'$,
transposing them changes the objective by $(p_{\sigma(r')} -
p_{\sigma(r)})(1/r - 1/r') > 0$; absence of improving transpositions is
exactly the descending sort --- the rearrangement inequality. The NDCG
claim is the same argument with $1/r$ replaced by the discount sequence.
Under multiple deals per window the argument applies to each successive
deal with the PL-updated distribution; ranking by the current $p_i$
remains optimal for the first hit, which is what recall@$K$ scores.
\end{proof}

The lemma settles the division of labor. A correctly specified,
likelihood-trained conditional model followed by a sort cannot be beaten
in expectation on recall@$K$, MRR, or NDCG; training on a surrogate can
help only by trading bias in regions the likelihood cares about for
accuracy in the top of the list when the score family cannot represent
the truth. That trade is real (B1), but it has a cost we refuse to pay in
the primary model: a LambdaMART scorer emits uncalibrated ordinal scores,
while the decision layer of \cref{ch:decision} consumes \emph{probabilities}
--- conditional identity probabilities composed with ground intensities
to produce statements like ``firm $i$ is on the list \emph{and} a deal is
likely within the quarter'' (T1 composed with T2/T3), and the conditional
NLL is the strictly proper score for that object
\citep{gneitingraftery2007}. A scorer that cannot feed the composition is
a dead end in this architecture regardless of its list metrics.

\begin{decision}{D10.3 --- LambdaMART as challenger only}
LambdaMART (and pairwise RankNet as its diagnostic little sibling) is
maintained as a benchmarked \emph{challenger} in the evaluation harness,
never as the primary model. Rationale: \cref{lem:ch10:topk} says it can
only win through score-family misspecification, which the boosted
grouped-softmax primary already addresses inside the likelihood family;
and its output does not feed the probability-consuming decision layer.
Reversal trigger: if LambdaMART beats the likelihood-trained boosted
ranker on pooled recall@5 by more than the era-blocked bootstrap interval
for two consecutive campaigns, \emph{and} a post-hoc calibration overlay
\citep{platt1999,zadroznyelkan2002,niculescumizil2005} on its scores
passes the calibration checks of \cref{ch:gof} at the event-group level,
it is promoted to primary and this section is rewritten around the
surrogate-first architecture.
\end{decision}

\section{The recall at $K$ protocol}\label{ch10:sec:protocol}

Everything above concerns fitting; what follows is the committed
evaluation contract --- the repository's headline metric, stated fully so
that no future implementation ``simplifies'' it into something
unauditable.

\begin{protocol}[Rolling recall@$K$ and MRR for the mark factor]\label{prot:ch10:recall}
\leavevmode
\begin{enumerate}[label=(\alph*)]
\item \emph{Origins.} Evaluation origins $\tau$ run on a weekly grid over
the held-out era. At each $\tau$ and each sector $s$, freeze the risk set
$\riskset_s(\tau)$ under the versioned risk-set definition (item (g)),
compute point-in-time scores, and emit the full ordering of
$\riskset_s(\tau)$; the top-$K$ prefixes are the lists $L_{s,\tau,K}$,
$K \in \{1, 5, 10\}$.
\item \emph{Window and hits.} The evaluation window is the next $W$
weeks, $(\tau, \tau + W]$; we commit to next-$W$-weeks rather than
next-$K'$-deals so that empty windows are informative and the metric has
a fixed temporal meaning. $W$ is matched to the consuming action's
refresh horizon: default $W = 4$ (a monthly sourcing look-ahead), with $W
= 13$ reported as a robustness column. A deal in the window scores a
\emph{hit} at level $K$ iff its firm was on $L_{s,\tau,K}$ --- membership
at $\tau$, not at deal time.
\item \emph{Denominator and newcomers.} The denominator counts deals in
the window whose firm was in $\riskset_s(\tau)$. Deals to firms outside
every frozen risk set at $\tau$ --- first appearances and out-of-universe
deals --- are \emph{excluded from the denominator and reported
separately} as the newcomer share, under the newcomer-vs-entrant
discipline of \cref{ch:universe}. Silently scoring newcomer deals as
misses (or dropping them without reporting) is forbidden: the first
deflates the metric for a failure the mark factor cannot commit, the
second hides the (C6) exposure.
\item \emph{Metrics.} Report recall@$K$ for $K \in \{1,5,10\}$ and MRR
(reciprocal rank of the funded firm in the full frozen ordering),
per sector and pooled (deal-weighted), aggregated over origins
rolling-origin style. Alongside, \emph{always}, the held-out conditional
NLL of \cref{eq:ch10:nll} on the same events: the list metrics say
whether the sort is right, the NLL says whether the probabilities feeding
the decision layer are right, and \cref{lem:ch10:topk} is the license to
demand both from one model.
\item \emph{Uncertainty.} Era-blocked bootstrap: resample contiguous
blocks of evaluation origins (blocks aligned with the macro-era
segmentation of \cref{ch:gof}) with replacement, recompute pooled
metrics, report percentile intervals. Origin-level i.i.d.\ resampling is
forbidden --- overlapping windows and regime persistence make it
anti-conservative.
\item \emph{Baselines.} Every report carries the same metrics for the
random ranker over the same frozen risk sets (chance floor) and, on
synthetic data, the oracle ranker (ceiling); the rehearsal reference
values are in \cref{tab:ch10:rehearsal}.
\item \emph{Frozen, versioned risk sets.} The risk-set definition ---
entry rule, exposure/exit rule from \cref{ch:buckets}, sector assignment
--- is frozen per evaluation campaign and identified by a version string
recorded with every metric row; changing the rule starts a new metric
series rather than overwriting one. Rationale: risk-set truthfulness is
the single largest sensitivity in the system --- rehearsal top-5 was
$0.14$ on observed (contaminated) risk sets against $0.29$ under oracle
membership (\cref{app:results}) --- so an unversioned membership rule
would make every metric movement uninterpretable.
\end{enumerate}
\end{protocol}

\begin{table}[htbp]
\centering
\small
\begin{tabular}{lccc}
\toprule
\textbf{Held-out metric} & \textbf{fitted} & \textbf{oracle (ceiling)} &
\textbf{random (floor)} \\
\midrule
top-1 hit rate & 16.7\% & 21.2\% & 2.9\% \\
top-5 hit rate & 49.2\% & 53.7\% & 14.3\% \\
MRR            & 0.328  & 0.371  & 0.118 \\
\bottomrule
\end{tabular}
\caption{Rehearsal reference numbers for the mark-factor ranker
(\cref{app:results}): the fitted ranker recovers most of the oracle's
performance, and the oracle itself is far from $1$ --- the residual is
irreducible softmax stochasticity, not model error. These values
calibrate expectations for real-data campaigns; they are priors and
vetoes, never substitutes (\cref{ch02:sec:synthetic}).}
\label{tab:ch10:rehearsal}
\end{table}

\begin{decision}{D10.1 --- the T1 deliverable and its two metrics}
The primary T1 deliverable is the sector-conditional risk-set ranker:
gradient-boosted scores inside the grouped-softmax conditional likelihood
\cref{eq:ch10:nll} (per \cref{lem:ch10:gradhess}), with cooldown and
staleness-age features, sampled risk sets per \cref{thm:ch10:sampled}
where scale requires, upgraded to Plackett--Luce per D10.2. It is
evaluated by \cref{prot:ch10:recall} \emph{and} held-out conditional NLL
--- both, always; a model improving one while degrading the other is
investigated, not shipped. Rationale: the competing-intensities identity
makes the conditional model the exact T1 target
(\cref{prop:ch10:competing}); the cancellation and consistency results
make it the defect-robust factor (\cref{prop:ch10:cancel,prop:ch10:consistency});
the B1 verdict selects boosted over linear scores within the family; and
\cref{lem:ch10:topk} guarantees likelihood training plus sorting is
list-optimal under the model while the NLL keeps the probabilities
honest for \cref{ch:decision}. Reversal triggers: if the boosted scorer
fails to beat the linear conditional logit on real data (both metrics,
two campaigns), demote to the linear score for interpretability; if the
challenger clause of D10.3 fires, the primary itself is reconsidered.
\end{decision}

\section{Defect declaration}\label{ch10:sec:defects}

Per design decision D2.1, the mark factor's position against the ledger
of \cref{tab:ch02:defects}, with the proofs above doing the work.

\emph{(C1) Administrative censoring: robust.} The conditional likelihood
\cref{eq:ch10:nll} is a sum over observed events only; no survival factor
appears because timing is conditioned away (\cref{prop:ch10:cancel}), and
under independent censoring the risk set at each observed event is
uncorrupted. Open spells cost the ranker nothing --- there is simply no
term.

\emph{(C2) Left truncation: robust for the conditional estimand,
with a scope caveat.} Comparisons are made within risk sets of firms that
have entered observation; the estimand is conditional on membership, so
delayed entry does not bias it. What (C2) forbids is re-reading scores as
statements about the untracked population --- that extension belongs to
\cref{ch:universe}.

\emph{(C3) Unlabeled exits: vulnerable --- the binding defect.} Dead
firms retained in $\riskset_s(t)$ occupy denominator mass and top-$K$
slots; \cref{prop:ch10:consistency}'s proof shows the estimand silently
becomes the conditional law given the contaminated pool. The rehearsal
quantified the cost as the largest sensitivity in the whole system
(top-5 $0.14$ observed vs $0.29$ oracle; an exit classifier with AUC
$0.877$ was not sufficient to close the gap, \cref{app:results}).
Mitigations are the exposure rules of \cref{ch:buckets}
(trailing-activity windows closing spells point-in-time) plus staleness
features that let the scorer learn to demote ghost-like profiles. One
inherited caution must be stated: the rehearsal's inverted-weight result
shows that the Bayes-optimal score weighting \emph{given a contaminated
pool} differs from the oracle weighting --- the fitted model partially
adapts to contamination, so cleaning the risk set changes the optimal
parameters, and oracle-world or cleaned-world coefficients must never be
transplanted onto contaminated risk sets or vice versa
(\cref{app:results}). Risk-set definition and fitted model are a matched
pair; \cref{prot:ch10:recall}(g) versions them together.

\emph{(C4) Untrusted dates and aggregation: robust, by proof.}
\Cref{prop:ch10:cancel}: binning and jitter deform $\lam_0(s,t)$, which
cancels; the residual exposure is only through snapshot assignment
(a deal jittered across a week boundary meets a one-week-shifted risk set
and covariates) and through within-cell order, which D10.2's
exchangeable-order Plackett--Luce removes.

\emph{(C5) Stale covariates: vulnerable, with signed bias, remedies
inherited.} Stale $\zfeat_i$ is classical covariate measurement error
inside a choice model: score coefficients attenuate toward zero, the
softmax flattens, lists diffuse, and recall@$K$ degrades toward the
random floor --- a bias against the model, never in its favor. The
remedies of \cref{ch:stale} (staleness ages as first-class features,
covariate-state layer where it pays) apply verbatim to the ranker's
feature block, and the boosted scorer can learn staleness-dependent
discounting that the linear score cannot.

\emph{(C6) Out-of-universe deals: vulnerable without the outside option.}
Deals from untracked firms mean the true choice set exceeds
$\riskset_s(t)$; a mark factor without an outside option overstates every
tracked firm's probability by the factor $1/(1 - \pi_{\mathrm{new}})$.
The treatment --- the newcomer branch with its exact separation of binary
and choice likelihoods --- is \cref{ch:universe}'s; this chapter's
protocol already quarantines the exposure by excluding newcomer deals
from the recall denominator and reporting their share
(\cref{prot:ch10:recall}(c)).

\begin{codehooks}
Existing. \modrepo{sector\_ranker.py}: \code{fit\_startup\_ranker}
implements \cref{eq:ch10:score,eq:ch10:softmax,eq:ch10:nll} with sector
deviations, the $\eta_s \le 0$ bound of \cref{ch10:sec:cooldown}
(\code{constrain\_cooldown\_negative}), and dynamic risk sets via
\code{candidate\_set} and \code{cooldown\_vector};
\code{ranker\_predict\_proba} exposes the conditional probabilities;
\code{evaluate\_ranker} computes held-out conditional NLL, MRR, and
top-$k$ hit rates and is the seed of the harness below.
\modrepo{models/hawkes\_ranker.py} carries the regime-B rehearsal ranker
behind the cooldown result of \cref{ch10:sec:cooldown}.

To be written. (i) The grouped-softmax XGBoost objective module
(\modrepo{ranking/grouped\_softmax\_objective.py}): per-group softmax
from margin scores, per-candidate $(g_j, h_j)$ of
\cref{eq:ch10:gradhess}, group boundaries from event ids, uniform
risk-set sampling per \cref{thm:ch10:sampled} with the winner retained.
(ii) The Plackett--Luce upgrade of D10.2 inside
\modrepo{sector\_ranker.py} (sequential denominators, exchangeable-order
averaging for $d \le 3$). (iii) The evaluation harness
(\modrepo{eval/ranking.py}) implementing \cref{prot:ch10:recall}
verbatim: weekly origins, frozen versioned risk sets
(\code{riskset\_version} recorded per metric row), $W$-week windows,
newcomer-share reporting, era-blocked bootstrap, and the
random/oracle baseline columns of \cref{tab:ch10:rehearsal}. (iv) The
LambdaMART challenger of D10.3 as a harness plug-in, never as a
dependency of the decision layer.
\end{codehooks}
